\begin{frame}{Operadores lógicos: combinar condições}
\begin{center}
\begin{tabular}{lll}
\toprule
Operador & Nome & Resultado verdadeiro quando \\
\midrule
\texttt{\&\&} & E lógico & Ambas são verdadeiras \\
\texttt{||} & OU lógico & Ao menos uma é verdadeira \\
\texttt{!} & NÃO lógico & A condição é falsa \\
\bottomrule
\end{tabular}
\end{center}
\begin{itemize}
\item Em um contexto lógico, zero é falso; qualquer valor diferente de zero é verdadeiro.
\item \texttt{\&\&}, \texttt{||} e \texttt{!} produzem 0 ou 1, com tipo \texttt{int}.
\item Prefira condições nomeadas com \texttt{bool} para expressar estados.
\end{itemize}
\end{frame}
